\section{Nivel 2 --- Control regulatorio}

\begin{frame}{Nivel 2: objetivo}
  \begin{itemize}
    \item Genera los comandos de \textbf{torque} $T^\ast$ a partir de referencias de
          \textbf{velocidad} del supervisor.
    \item Control en cascada: \textbf{velocidad} (PI) $\to$ modulador de torque.
    \item Atenúa perturbaciones: balanceo de la carga y variaciones de masa.
    \item Opera muestreado en tiempo discreto.
  \end{itemize}
  \begin{block}{Componentes diseñados}
    Lazo PI de velocidad \quad$\bullet$\quad Observador de estado \quad$\bullet$\quad
    Control anti-balanceo \quad$\bullet$\quad Anti-windup \quad$\bullet$\quad
    Compensación de gravedad
  \end{block}
\end{frame}

\begin{frame}{Lazo de velocidad: PI por asignación de polos}
  \begin{columns}[T]
    \begin{column}{0.5\textwidth}
      Planta simplificada del carro:
      \[ G_{vel,t}(s)=\frac{V_t(s)}{T_{tm}(s)}=\frac{K_f}{M_t\,s+b_t}\approx\frac{K_f}{M_t\,s} \]
      Controlador:
      \[ C_{vel}(s)=K_p+\frac{K_i}{s} \]
      Comparando con $s^2+2\zeta\omega_n s+\omega_n^2$:
      \[ K_p=\frac{2\zeta\omega_n M_t}{K_f},\quad K_i=\frac{\omega_n^2 M_t}{K_f} \]
      \footnotesize $\omega_n=0.8$ rad/s, $\zeta=1$ (sin sobrepaso).
    \end{column}
    \begin{column}{0.5\textwidth}
      \centering
      \includegraphics[width=\linewidth]{nivel_2/lazo_controlador_carro.png}\\[0.2cm]
      \footnotesize \textbf{Anti-windup}: la señal de freno resetea el integrador mientras el
      freno está cerrado.
    \end{column}
  \end{columns}
\end{frame}

\begin{frame}{Control anti-balanceo (\textit{anti-sway})}
  \begin{columns}[T]
    \begin{column}{0.55\textwidth}
      Estrategia \textbf{desacoplada}: la velocidad del carro inyecta amortiguamiento sobre el
      péndulo.
      \[ v_{sway}(t)=K_{sway}\,\theta_l(t) \]
      Sintonía por linealización local ($\omega_n=\sqrt{g/l_h}$):
      \[ \boxed{\,K_{sway}(l_h)=2\sqrt{g\,l_h}\,} \]
      \begin{itemize}
        \item \textbf{Gain scheduling} con la altura de izaje $l_h$.
        \item Válido sin aproximación de ángulos pequeños ($\cos\theta_l>0$).
        \item Se subamortigua ($\zeta=0.5$) para no sobrecorregir la posición.
      \end{itemize}
    \end{column}
    \begin{column}{0.45\textwidth}
      \centering
      \includegraphics[width=\linewidth]{nivel_2/control_balanceo.png}
    \end{column}
  \end{columns}
\end{frame}

\begin{frame}{Observador de estado con acción integral}
  \begin{columns}[T]
    \begin{column}{0.5\textwidth}
      \begin{itemize}
        \item Solo se mide \textbf{posición} (encoder); el PI necesita la \textbf{velocidad}.
        \item Observador que estima $\tilde\omega$ a partir de $\theta$.
        \item \textbf{Acción integral} sobre el error de observación $\Rightarrow$ elimina el
              error en régimen ante perturbaciones constantes (balanceo).
        \item Sintonía por separación de anchos de banda: $20\times$ más rápido que la planta.
      \end{itemize}
    \end{column}
    \begin{column}{0.5\textwidth}
      \centering
      \includegraphics[width=\linewidth]{nivel_2/obs_pid_carro.png}
    \end{column}
  \end{columns}
\end{frame}

\begin{frame}{Izaje: gravedad y límite de potencia}
  \begin{columns}[T]
    \begin{column}{0.5\textwidth}
      \begin{block}{Compensación de gravedad (feedforward)}
        \[ T_{g,ff}=\frac{m_l\,g\,r_{hd}}{2\,i_h} \]
        El PI sólo corrige errores de modelo; el peso estático lo sostiene el feedforward.
      \end{block}
      \begin{block}{Operación a potencia constante}
        El límite de velocidad de izaje se reduce con la masa estimada para no exceder la
        potencia nominal.
      \end{block}
    \end{column}
    \begin{column}{0.5\textwidth}
      \centering
      \includegraphics[width=\linewidth]{nivel_2/lazo_controlador_izaje.png}
    \end{column}
  \end{columns}
\end{frame}
